\section{Twisted Goren--Oort Tanner codes}\label{sec:arith}

This section is not formalized. It supplies the algebraic half of Theorem T: an equivariant datum with MDS local codes
on the Bruhat--Tits tree of $\PGL_2(\Q_{p^2})$, and a tower of congruence quotients whose Tanner codes have rate close to
the square of the local rate (\cref{thm:G}). The twelfth paper of the series proves the case $p=3$ in weight $2$ in
detail~\cite[Section~6]{Luck2026}, from the same sources, and we follow it. Three steps are new: the Riemann--Roch
estimate in weight~$k$, where the Hasse invariant gives the vanishing (\cref{lem:RRk}); the girth bound
(\cref{lem:girth}); and rigidity (\cref{lem:rigid}), which uses the group $\SL_2(\Q_{p^2})$ rather than the icosian group,
because for $p>3$ the icosian group has many orbits on the edges of the tree.

\subsection{The icosian group at an inert prime}\label{sec:arith-group}

Let $E=\Q(\sqrt5)$, with ring of integers $\OO_E=\Z[\theta]$, $\theta=(1+\sqrt5)/2$; it has narrow class number one,
$\theta$ has norm $-1$, and every totally positive unit is a square. Let $B$ be the quaternion algebra over $E$ ramified
exactly at the two infinite places, and $I\subset B$ the icosian ring, a maximal order of class number one and type
number one whose norm-one units form the binary icosahedral group $2I\cong\SL_2(\F_5)$ of
order~$120$~\cite[Proposition~V.3.4, Corollaire~I.4.11 and p.~155]{Vigneras1980}.

Fix a prime $p\ge7$ inert in $E$, that is, $p\equiv\pm2\pmod5$. The completion of $E$ at $p$ is the unramified quadratic
extension $\Q_{p^2}$ of $\Q_p$, with ring of integers $\Z_{p^2}$ and residue field $\F_{p^2}$, and $B$ splits there:
$B\otimes_E\Q_{p^2}\cong M_2(\Q_{p^2})$ with $I\otimes\Z_{p^2}\cong M_2(\Z_{p^2})$. The group
\[
  \Gamma^+=I[1/p]^1/\{\pm1\}\subset\PSL_2(\Q_{p^2})
\]
acts on the Bruhat--Tits tree $\Tree$ of $\PGL_2(\Q_{p^2})$. Its vertices are the homothety classes of
$\Z_{p^2}$-lattices in $\Q_{p^2}^2$; the neighbours of $[\Lambda]$ are the classes of the lattices $\Lambda'$ with
$p\Lambda\subset\Lambda'\subset\Lambda$ and $\Lambda/\Lambda'\cong\F_{p^2}$, so the edges at $[\Lambda]$ correspond to
the points of $\Pone(\Lambda/p\Lambda)$, and $\Tree$ is $(p^2+1)$-regular. The vertices have two parities, the parity of
$\log_p\abs{\det}$ of a basis, which elements of determinant one preserve. Let $\Lambda_0$ be the vertex of
$I\otimes\Z_{p^2}$.

\begin{lemma}[The icosian group and its congruence subgroups]\label{lem:group}
\notformalized
\begin{enumerate}[label=(\alph*)]
\item There is $\alpha_0\in I$ of reduced norm $p$, and $\Lambda_1=\alpha_0\Lambda_0$ is a neighbour of $\Lambda_0$.
The group $\Gamma^+$ acts on $\Tree$ without inversions and transitively on the vertices of each parity; the stabilizers
of $\Lambda_0$ and $\Lambda_1$ are $A_5=2I/\{\pm1\}$ and $\alpha_0A_5\alpha_0^{-1}$.
\item Let $\lambda_0$ be a prime of $\OO_E$ of prime norm $\ell\ge7$, and for $j\ge1$ let
$\Gamma^+(\lambda_0^j)=\ker\bigl(\Gamma^+\to\PSL_2(\OO_E/\lambda_0^j)\bigr)$. The map is onto,
$\Gamma^+(\lambda_0^j)$ is torsion-free and acts freely on $\Tree$, the groups $\Gamma^+(\lambda_0^j)$ strictly decrease
with trivial intersection, and $X_j=\Gamma^+(\lambda_0^j)\backslash\Tree$ has
\[
  c_j=\abs{\PSL_2(\OO_E/\lambda_0^j)}/60
\]
vertices of each parity.
\end{enumerate}
\end{lemma}
\begin{informalproof}
(a) This is the argument of~\cite[Lemma~6.2]{Luck2026} with $3$ replaced by $p$. A vertex is the class of the completion
at $p$ of a right $I$-ideal $\mathfrak a$ with $\mathfrak a_{\mathfrak q}=I_{\mathfrak q}$ for $\mathfrak q\nmid p$, and
$\mathfrak a=\beta I$ with $\beta\in I[1/p]^\times$ because $I$ has class number one. The reduced norm of $\beta$ is a
totally positive unit of $\OO_E[1/p]$, hence $p^a\theta^{2b}$, and we may take $b=0$. Applied to a neighbour $\Lambda'$ of
$\Lambda_0$ with $\Lambda'\subset\Lambda_0$, whose ideal has reduced norm the prime ideal $(p)$, this gives $\alpha_0$.
For a general vertex, $\beta/p^{a/2}$ (for even $a$) or $(\beta/p^{(a-1)/2})\alpha_0^{-1}$ (for odd $a$) lies in
$I[1/p]^1$ and moves $\Lambda_0$ or $\Lambda_1$ to it. Elements of reduced norm one preserve parity, so there are no
inversions. The stabilizer of $\Lambda_0$ in $I[1/p]^1$ consists of the elements that are integral at $p$, which is
$I^1=2I$; conjugation by $\alpha_0\in I[1/p]^\times$ preserves $I[1/p]^1$ and gives the stabilizer of $\Lambda_1$.

(b) This is~\cite[Lemma~6.3]{Luck2026} with $3$ and $11$ replaced by $p$ and $\ell$. Strong approximation for the
simply connected group $B^1$, which is not compact at $p$~\cite[Section~7.4]{PlatonovRapinchuk1994}, gives
surjectivity. A nontrivial element of finite order of $\Gamma^+$ fixes a vertex~\cite[Section~I.4]{Serre1980}, so it is
conjugate into $A_5$ or $\alpha_0A_5\alpha_0^{-1}$, which map injectively to $\PSL_2(\OO_E/\lambda_0)$ because
$1+\lambda_0M_2(\OO_{E,\lambda_0})$ is torsion-free for odd~$\ell$. Hence $\Gamma^+(\lambda_0^j)$ is torsion-free and,
acting without inversions, acts freely. It is normal and meets no vertex stabilizer, so the vertices of $X_j$ of a given
parity form the double coset space $\Gamma^+(\lambda_0^j)\backslash\Gamma^+/A_5$, of size
$\abs{\PSL_2(\OO_E/\lambda_0^j)}/60$, and likewise for the other parity. The decrease and the trivial intersection are
proved as in~\cite[Lemma~6.3]{Luck2026}.
\end{informalproof}

So $X_j$ is bipartite, with $c_j$ vertices of each parity and $(p^2+1)c_j$ edges, and it is $(p^2+1)$-regular because
$s(\Gamma^+(\lambda_0^j))\ge2$ (\cref{lem:girth} below, with $r=0$, and \cref{sec:prelim-graphs}); it is the dual graph
of the Goren--Oort divisor below. Its girth grows along the tower, by an archimedean argument in the style of
Lubotzky, Phillips and Sarnak~\cite{LPS1988}.

\begin{lemma}[Girth]\label{lem:girth}
\notformalized
Let $j\ge1$ and $r\ge0$ with $4p^r<\ell^j$. Then every element of $\Gamma^+(\lambda_0^j)\setminus\{1\}$ moves every
vertex of $\Tree$ by at least $2r+2$. Hence $s(\Gamma^+(\lambda_0^j))\ge2r+2$; in particular $X_1$ is simple when
$\ell>4p$, and $s(\Gamma^+(\lambda_0^j))\ge2(j\log\ell-\log4)/\log p$ for all~$j$.
\end{lemma}
\begin{informalproof}
Elements of $\Gamma^+$ preserve parity, so they move vertices by even distances. Suppose $\gamma\ne1$ in
$\Gamma^+(\lambda_0^j)$ moves a vertex $x$ by $2r'$ with $1\le r'\le r$. By \cref{lem:group}(a), $x=g\Lambda_0$ or
$x=g\alpha_0\Lambda_0$ with $g\in\Gamma^+$; conjugating by $g$ and by $\alpha_0$, which preserve $\Gamma^+(\lambda_0^j)$
(the first because it is normal, the second because $\alpha_0$ is a unit at $\lambda_0$), we may assume $x=\Lambda_0$.
Let $\tilde\gamma\in I[1/p]^1$ be the lift with $\tilde\gamma\equiv1\pmod{\lambda_0^j}$; it is unique because
$-1\not\equiv1\pmod{\lambda_0}$~\cite[Lemma~6.3]{Luck2026}. By the Cartan decomposition,
$\alpha=p^{r'}\tilde\gamma$ lies in $M_2(\Z_{p^2})=I\otimes\Z_{p^2}$, hence in $I$, and $\alpha\equiv p^{r'}\pmod{\lambda_0^jI}$.
Let $\pi$ be a totally positive generator of $\lambda_0^j$ and write $\alpha=p^{r'}+\pi\beta$ with $\beta\in I$. Then
$\beta\ne0$ because $\gamma\ne1$, and
\[
  p^{2r'}=\operatorname{Nrd}\alpha=p^{2r'}+p^{r'}\pi\operatorname{Tr}\beta+\pi^2\operatorname{Nrd}\beta
\]
gives $p^{r'}\operatorname{Tr}\beta=-\pi\operatorname{Nrd}\beta$. Since $\pi$ is prime to $p$, $\operatorname{Tr}\beta=\pi u$
with $u\in\OO_E$, and then $\operatorname{Nrd}\beta=-p^{r'}u$, so $u\ne0$. At each real embedding $\iota$ of $E$,
$B\otimes_\iota\R$ is the algebra of Hamilton quaternions, so
$\iota(\operatorname{Tr}\beta)^2\le4\,\iota(\operatorname{Nrd}\beta)$, that is,
$\iota(\pi)^2\iota(u)^2\le4p^{r'}\abs{\iota(u)}$, hence $\iota(\pi)^2\abs{\iota(u)}\le4p^{r'}$. Multiplying over the two
embeddings gives $\ell^{2j}\abs{N(u)}\le16p^{2r'}$, and $\abs{N(u)}\ge1$, so $\ell^j\le4p^{r'}\le4p^r$, a contradiction.
The displacement bound follows from~\eqref{eq:disp}. For $j=1$ and $r=1$ the quotient is simple by
\cref{sec:prelim-graphs}. The last bound takes the largest admissible~$r$.
\end{informalproof}

\subsection{Riemann--Roch in weight \texorpdfstring{$k$}{k}}\label{sec:arith-RR}

Let $\mathfrak n=\lambda_0^j$, an ideal prime to $p$ with a totally positive generator, and
$\mathcal K(\mathfrak n)\subset\GL_2(\hat\OO_E)$ its principal congruence subgroup. It is neat: a root of unity congruent
to $1$ modulo a prime above $\ell$ has $\ell$-power order, so it is $1$ or has degree at least $\ell-1\ge6$ over $\Q$,
while eigenvalues of elements of $\GL_2(E)$ have degree at most~$4$. As in~\cite[Section~6.2]{Luck2026}, choose a smooth
projective toroidal compactification over $W(\bar\F_p)$ of the closure of one geometric connected component of the
Hilbert modular surface of $E$ of level $\mathcal K(\mathfrak n)$, and let $S$ be its special fibre, with boundary $D$.
We take the one given at every cusp by Hirzebruch's minimal cone decomposition~\cite[Section~2.5]{Hirzebruch1973}. Every
cusp has type $(\OO_E,\mathcal U_{\mathfrak n})$ up to a totally positive scalar~\cite[Section~6.2, (G4)]{Luck2026}, and
$\theta^2=3-\theta^{-2}$ has the purely periodic continued fraction $[[3]]$, so the rays of the decomposition are spanned by
lattice points $A_k$ ($k\in\Z$) with $A_{k-1}+A_{k+1}=3A_k$, any two consecutive ones forming a basis of the
lattice~\cite[Sections~2.5 and~2.6]{Hirzebruch1973}. So the decomposition is smooth. It is also projective in the sense
of~\cite[Definitions~7.3.1.1 and~7.3.1.3]{Lan2013}. Let $\varphi$ be the function that is linear on each cone and equal
to $1$ at every $A_k$. It takes positive integer values at nonzero lattice points, by smoothness. If
$x=aA_{k-1}+a'A_k$ and $y=c'A_k+cA_{k+1}$ with nonnegative coefficients, then
$\varphi(x+y)-\varphi(x)-\varphi(y)=\min(a,c)$, by $A_{k-1}+A_{k+1}=3A_k$. So $\varphi$ bends at every ray, hence it is
concave (superadditive, Lan's convexity) and its maximal domains of linearity are exactly the cones: it is a
polarization function. The complex fibre of the component is $\Gamma(\mathfrak n)\backslash\mathfrak H^2$ with
$\Gamma(\mathfrak n)\subset\SL_2(\OO_E)$. We use the facts (G1)--(G4) of~\cite[Section~6.2]{Luck2026}, whose sources hold
for every prime unramified in $E$ and every level prime to it; we restate them for~$p$.
\begin{itemize}
\item[(G1)] $S$ is smooth, projective and connected, and $D\ne0$ is a disjoint union of cycles of smooth rational curves,
one for each cusp~\cite[Theorem~5.1, Proposition~5.2 and Corollary~5.3]{Rapoport1978},
\cite{Chai1990}, \cite[Section~7]{Dimitrov2004}, \cite[Theorems~6.4.1.1 and~7.3.3.4]{Lan2013}.
\item[(G2)] The Hodge bundle splits as $\omega_1\oplus\omega_2$ because $p$ is unramified, and
$K_S+D\cong\omega_1^2\otimes\omega_2^2$~\cite[Theorem~6.4.1.1(4)]{Lan2013}. Write $H=\omega_1\otimes\omega_2$ and, for an
integer $k$, $L_k=H^{\otimes k}$, the bundle of parallel weight~$k$; so $K_S+D=2H$ in additive notation.
\item[(G3)] The partial Hasse invariants $h_1\in H^0(\omega_1^{-1}\omega_2^p)$ and $h_2\in H^0(\omega_1^p\omega_2^{-1})$
have reduced divisors $Z_1$, $Z_2$, regular curves meeting in finitely many
points~\cite[Section~1 and Theorem~2.1(1)]{Goren2001}; they extend to $S$ by the Koecher
principle~\cite[Theorem~8.3]{Dimitrov2004} and do not vanish on $D$, because their $q$-expansions at every cusp are
$1$~\cite[Theorem~2.1(2)]{Goren2001}. So $h=h_1h_2\in H^0(L_{p-1})$ has divisor $Y=Z_1+Z_2$, and $\OO(Y)\cong L_{p-1}$.
Since $p$ is inert, $Y$ is the reduced supersingular locus~\cite[Proposition~4.1]{BachmatGoren1999}.
\item[(G4)] Let $e=\abs{\PSL_2(\OO_E/\mathfrak n)}\cdot2\zeta_E(-1)=\abs{\PSL_2(\OO_E/\mathfrak n)}/15$, the Euler
number of $\Gamma(\mathfrak n)\backslash\mathfrak H^2$~\cite[Sections~1.6 and~1.7]{Hirzebruch1973}. Every cusp is resolved
by a cycle of curves of self-intersection $-3$~\cite[Sections~2.5 and~2.6]{Hirzebruch1973}, and
$K_S^2=2e+\sum(2-b)=2e-\#\{\text{boundary curves}\}$, where the sum runs over the boundary curves, of self-intersection
$-b=-3$~\cite[Sections~3.2 and~3.6]{Hirzebruch1973}. Noether's formula $\chi(\OO_S)=(K_S^2+c_2(S))/12$, with
$c_2(S)=e+\#\{\text{boundary curves}\}$, then gives $\chi(\OO_S)=e/4=c_j$, as in~\cite[Section~6.2, (G4)]{Luck2026}. The
boundary curves over each cusp are the closures of the strata of the rays of the decomposition, which are smooth over
$W(\bar\F_p)$~\cite[Theorem~6.4.1.1(2) and~(3)]{Lan2013}, so the configuration of the boundary is the same on $S$ and on the
complex fibre; $\chi(\OO)$ and intersection numbers of line bundles are constant in the smooth proper family over
$W(\bar\F_p)$. So these facts may be checked on the complex fibre, where they are Hirzebruch's. Each cycle has at least five
curves: its length is the index of the squares of the units congruent to $1$ modulo $\mathfrak n$ in the totally positive
units~\cite[Section~6.2, (G4)]{Luck2026}, which is the order of $\theta$ in $(\OO_E/\mathfrak n)^\times/\{\pm1\}$, and for
$1\le a\le4$ the numbers $\theta^a\mp1$ have norms $\pm1$, $\pm4$, $\pm5$ or~$9$, which no prime $\ell\ge7$ divides.
\end{itemize}

\begin{lemma}[Sections in weight $k$]\label{lem:RRk}
\notformalized
For $2\le k\le p-2$, $h^0(S,L_k)\ge(k-1)^2c_j$, and restriction $H^0(S,L_k)\to H^0(Y,L_k|_Y)$ is injective. Hence
$h^0(Y,L_k|_Y)\ge(k-1)^2c_j$.
\end{lemma}
\begin{informalproof}
\emph{Intersection numbers.} Each boundary cycle $D_i$ has arithmetic genus $1$, and the cycles are disjoint, so
adjunction gives $(K_S+D)\cdot D=0$; hence $H\cdot D=0$ and $K_S\cdot D=-D^2$. A cycle of curves of
self-intersections $-b$ has square $\sum(2-b)$, so $D^2=-\#\{\text{boundary curves}\}$. Then
$(2H)^2=(K_S+D)^2=K_S^2-D^2=2e$ by (G4), so $H^2=e/2=2c_j$, and $H\cdot K_S=H\cdot(2H-D)=4c_j$.

\emph{Riemann--Roch.} For a line bundle $L$ on the smooth projective surface $S$,
$\chi(L)=\chi(\OO_S)+\frac12L\cdot(L-K_S)$~\cite[Section~4.1]{Liedtke2013}. For $L=L_k=kH$ this is
\[
  \chi(L_k)=c_j+\tfrac12\bigl(k^2H^2-kH\cdot K_S\bigr)=c_j+k^2c_j-2kc_j=(k-1)^2c_j.
\]

\emph{Vanishing.} By Serre duality $h^2(L_k)=h^0(K_S-L_k)=h^0((2-k)H-D)$. A nonzero section of this bundle, multiplied
by a nonzero section of $\OO(D)$, gives a nonzero section of $-(k-2)H$ when $k\ge3$, and is a nonzero section of $\OO(-D)$
when $k=2$. The latter does not exist because $D$ is a nonzero effective divisor on the integral projective surface $S$.
If $s\ne0$ is a section of $-aH$ with $a\ge1$, then $s^{p-1}h^a$ is a nonzero section of $\OO_S$, hence a nonzero
constant, so $h$ vanishes nowhere and $(p-1)H\cong\OO(Y)$ is trivial by (G3). This contradicts $H^2=2c_j>0$. So
$h^2(L_k)=0$ and $h^0(L_k)\ge\chi(L_k)=(k-1)^2c_j$.

\emph{Injectivity.} The kernel of restriction to $Y$ is $H^0(L_k(-Y))=H^0(-(p-1-k)H)$, by (G3), and $p-1-k\ge1$; it
vanishes by the same argument.
\end{informalproof}

\subsection{Uniformization of the supersingular locus}\label{sec:arith-unif}

Let $\mathcal M$ be the Rapoport--Zink space of supersingular polarized $p$-divisible groups with $\OO_E\otimes\Z_p$-action,
$J$ the group of $\OO_E$-linear self-quasi-isogenies of its framing object that multiply the polarization by a factor in
$\Q_p^\times$, and $\mathcal M^{(n)}_{\mathrm{red}}$ the part of its reduced
locus described in~\cite[Section~6.3]{Luck2026}, after Kudla and Rapoport~\cite[Section~4]{KudlaRapoport1999} and
Stamm~\cite{Stamm1997}. It is a union of projective lines: a line $C_\Lambda=\Pone(\Lambda/p\Lambda)$ of the first kind for
each $\Z_{p^2}$-lattice $\Lambda$ of a two-dimensional $\Q_{p^2}$-space $U_0$ with $\langle\Lambda,\Lambda\rangle=p^n\Z_{p^2}$,
and a line of the second kind for each such lattice of a second space $U_1$, which a $\sigma$-semilinear bijection
$\Phi\colon U_1\to U_0$ commuting with $J$ identifies with $U_0$. Here $\OO_E\otimes\Z_p$ acts on $U_0$ by scalars,
through the embedding into $W=W(\bar\F_p)$ that defines the grading of the isocrystal~\cite[(4.3)]{KudlaRapoport1999}.
Lines of the same kind are disjoint; each line meets
exactly one line of the other kind at each of its $p^2+1$ points rational over $\F_{p^2}$, the superspecial points, and no
line elsewhere~\cite[Lemma~4.3(iii), the text after~(4.9) and Proposition~4.4]{KudlaRapoport1999}; the crossings are ordinary double
points~\cite[Theorem~6.1]{BachmatGoren1999}. Sending a line of the first kind to $[\Lambda]$ and one of the second kind to
$[\Phi\Lambda']$ identifies the dual graph with the tree $\Tree_{U_0}$ of $U_0$, equivariantly for the subgroup $J^0$ of
elements whose similitude factor is a $p$-adic unit~\cite[Section~6.3]{Luck2026}. The Hodge bundles of the universal
$p$-divisible group, and so the bundles $L_k$, carry a canonical action of~$J$.

\begin{lemma}[Uniformization]\label{lem:unif}
\notformalized
There are an integer $n$, an action of $\Gamma^+$ on $\mathcal M^{(n)}_{\mathrm{red}}$ through $J^0$, and a bijection
$\Psi\colon U_0\to\Q_{p^2}^2$, either $\Q_{p^2}$-linear or $\sigma$-semilinear, carrying $\Tree_{U_0}$ to $\Tree$, such
that:
\begin{enumerate}[label=(\alph*)]
\item the identification of the dual graph of $\mathcal M^{(n)}_{\mathrm{red}}$ with $\Tree$ is $\Gamma^+$-equivariant;
\item $Y\cong\Gamma^+(\lambda_0^j)\backslash\mathcal M^{(n)}_{\mathrm{red}}$, and under this isomorphism $L_k|_Y$ is the
descent of $L_k$ with its $J$-action; in particular the dual graph of $Y$ is~$X_j$;
\item $L_k$ has degree $m=k(p-1)$ on every line of $\mathcal M^{(n)}_{\mathrm{red}}$.
\end{enumerate}
\end{lemma}
\begin{informalproof}
Parts (a) and (b) are~\cite[Lemma~6.6]{Luck2026} with $3$, $11$ and $\Gamma^+(\lambda_0^{k+1})$ replaced by $p$, $\ell$ and
$\Gamma^+(\lambda_0^j)$, and with $L_k$ in place of $L_2$; its proof uses the uniformization theorem of Rapoport and
Zink~\cite[Theorem~6.30]{RapoportZink1996}, the Hasse principle for the group of the moduli problem, Tate's
theorem~\cite{Tate1966}, the type number one of $I$, the norm theorem of Hasse and
Schilling~\cite[Theorem~III.4.1]{Vigneras1980} and strong approximation, none of which depends on $p$ or on the weight.
It also needs a supersingular point on $S$, which exists because $h^0(Y,L_k|_Y)\ge(k-1)^2c_j>0$ (\cref{lem:RRk}), so
$Y\ne\emptyset$. There the identification is chosen by the Skolem--Noether theorem: the action of $B$ on $U_0$, composed with the main
involution, becomes the isomorphism $B\otimes_E\Q_{p^2}\cong M_2(\Q_{p^2})$ of \cref{sec:arith-group}, composed with
$\sigma$ when the embedding of $\OO_E\otimes\Z_p$ into $W$ that defines the grading differs from the identification of
the completion of $E$ at $p$ with $\Q_{p^2}$ used there~\cite[proof of Lemma~6.6(a)]{Luck2026}. In that case we compose
the identification with $\sigma$ on coordinates. This makes $\Psi$ $\sigma$-semilinear and turns the action of $\Gamma^+$
on $\Q_{p^2}^2$, and so on $\Tree$, into the standard one.
The element $-1$ acts trivially on $\mathcal M^{(n)}_{\mathrm{red}}$ and on $L_k=(\omega_1\otimes\omega_2)^{\otimes k}$,
by $(-1)^{2k}=1$, so $I[1/p]^1$ acts through $\Gamma^+$.

(c) Since $\Gamma^+(\lambda_0^j)$ acts freely, every line maps isomorphically onto a component of $Y$, compatibly
with $L_k$. Let $C$ be a component of $Z_1$, and $a=\deg\omega_1|_C$, $b=\deg\omega_2|_C$. The section $h_2$ vanishes on
$C$ exactly at its $p^2+1$ superspecial points, simply, because the crossings are ordinary double points of the reduced
curve $Z_1+Z_2$; so $pa-b=p^2+1$. The components of $Z_1$ are disjoint and $Z_1\sim p\omega_2-\omega_1$ by (G3), so
$C^2=Z_1\cdot C=-a+pb$; and adjunction with $K_S=2H-D$ and $D\cdot C=0$ gives
$-2=(K_S+C)\cdot C=2a+2b-a+pb=a+(p+2)b$. Hence $(a,b)=(p,-1)$, and $\deg L_k|_C=k(a+b)=k(p-1)$. The components of $Z_2$ are
symmetric.
\end{informalproof}

A \emph{geometric datum} on $\Tree$ consists of a smooth rational curve $C_v$ over $\bar\F_p$ and a line bundle $L_v$ of
degree $m$ on it for every vertex $v$, a point $P_{v,e}\in C_v$ for every edge $e$ at $v$ (its \emph{special points}), and
for every edge $e=\{v,w\}$ an isomorphism $\varphi_e$ between the fibres of $L_v$ at $P_{v,e}$ and of $L_w$ at $P_{w,e}$.
For a group $\Gamma$ acting on the datum and freely on $\Tree$, its \emph{Tanner code} on $\Gamma\backslash\Tree$ is the space of
$\Gamma$-invariant families of sections $s_v\in H^0(C_v,L_v)$ whose values at the two special points of each edge
correspond under $\varphi_e$. By \cref{lem:unif}, the lines of $\mathcal M^{(n)}_{\mathrm{red}}$ with the restrictions of
$L_k$, their superspecial points and the identifications of the fibres of $L_k$ at the crossings form a geometric datum
$D_{\mathrm{geo}}$ on~$\Tree$, the same for every~$j$.

\begin{lemma}[Sections are Tanner codes]\label{lem:sections}
\notformalized
For every subgroup $\Gamma\subseteq\Gamma^+(\lambda_0)$ of finite index, $H^0(\Gamma\backslash\mathcal M^{(n)}_{\mathrm{red}},L_k)$
is isomorphic to the Tanner code of $D_{\mathrm{geo}}$ on $\Gamma\backslash\Tree$.
\end{lemma}
\begin{informalproof}
This is~\cite[Lemma~6.7]{Luck2026}: the curve $\Gamma\backslash\mathcal M^{(n)}_{\mathrm{red}}$ has the lines as
components and ordinary double points exactly at the superspecial points, so by the normalization sequence a section of
$L_k$ is a family of sections on the components whose values at each node correspond.
\end{informalproof}

\subsection{Rigidity and the twisted model over \texorpdfstring{$\F_{p^2}$}{F\_{p\textasciicircum2}}}\label{sec:arith-rigid}

Let $G=\SL_2(\Q_{p^2})$. It acts on $\Tree$ without inversions, transitively on the vertices of each parity and on the
edges. Let $v_0=[\Z_{p^2}^2]$, $v_1=[\Z_{p^2}e_1\oplus p\Z_{p^2}e_2]$ and $e_0=\{v_0,v_1\}$. The
stabilizers are $K_0=\SL_2(\Z_{p^2})$, $K_1=g_1K_0g_1^{-1}$ with $g_1=\operatorname{diag}(1,p)$, and the Iwahori subgroup
$K_0\cap K_1$ of matrices with lower left entry divisible by~$p$. Let $\rho_0\colon K_0\to\SL_2(\F_{p^2})$ be reduction modulo
$p$ and $\rho_1(g)=\rho_0(g_1^{-1}gg_1)$ on $K_1$; then $\rho_0$ maps the Iwahori subgroup onto the upper and $\rho_1$ onto
the lower triangular matrices, with the same diagonal: if $g=(\begin{smallmatrix}a&b\\pc&d\end{smallmatrix})$, then
$\rho_0(g)=(\begin{smallmatrix}\bar a&\bar b\\0&\bar d\end{smallmatrix})$ and
$\rho_1(g)=(\begin{smallmatrix}\bar a&0\\\bar c&\bar d\end{smallmatrix})$, with $\bar a\bar d=1$. Let $\sigma$ be the Frobenius
of $\F_{p^2}$, applied to matrices entrywise.

A \emph{$G$-equivariant} geometric datum carries an action of $G$ compatible with all its pieces, such that the stabilizer
$K$ of each vertex $v$ acts on $(C_v,L_v)$ through the reduction of $K$ modulo $p$, that is, through
$\SL(\Lambda_v/p\Lambda_v)\cong\SL_2(\F_{p^2})$ for a lattice $\Lambda_v$ of $v$. Its \emph{type} at $v_0$ (respectively
$v_1$) is \emph{standard} if $C_v$ with this action is isomorphic to $\Pone_{\bar\F_p}$ with the action through $\rho_0$
(respectively $\rho_1$), and \emph{twisted} if it is isomorphic to $\Pone_{\bar\F_p}$ with the action through
$\sigma\circ\rho_0$ (respectively $\sigma\circ\rho_1$).

\paragraph{The twisted model.} Let $m=k(p-1)$. At $v_0$ take $C_0=\Pone_{\F_{p^2}}$ with $K_0$ acting through $\rho_0$, at
$v_1$ take $C_1=\Pone_{\F_{p^2}}$ with $K_1$ acting through $\sigma\circ\rho_1$, and on both the bundle $\OO(m)$, whose
sections are the binary forms of degree $m$, with its natural linearization. The special points of $e_0$ are the points
$[e_1]\in C_0$ and $[e_2]\in C_1$: they are the unique fixed points of the images of the Iwahori subgroup, upper and lower
triangular. The fibre of $\OO(m)$ at $[e_1]$ has the basis $(e_1^\vee)^{\otimes m}$, on which the Iwahori subgroup acts by
$\bar a^{-m}$; the fibre at $[e_2]$ has the basis $(e_2^\vee)^{\otimes m}$, on which it acts by $\sigma(\bar d)^{-m}=\bar a^{pm}$.
Since $\bar a^{(p+1)m}=\bar a^{k(p^2-1)}=1$, the identification
$\varphi_{e_0}\colon(e_1^\vee)^{\otimes m}\mapsto(e_2^\vee)^{\otimes m}$ is equivariant. Transporting these data by~$G$, which
is possible because $G$ acts on $\Tree$ with the edge $e_0$ as a fundamental domain and every piece is equivariant for the
relevant stabilizer, defines a $G$-equivariant geometric datum $D_{\mathrm{tw}}$ of degree $m$, standard at the vertices
of the parity of $v_0$ and twisted at the others. Every piece of it is defined over $\F_{p^2}$, and $G$ acts through
$\F_{p^2}$-rational maps. The element $-1$ acts trivially, since $m$ is even, so subgroups of $\PSL_2(\Q_{p^2})$ act on it.
Let $D^\sigma_{\mathrm{tw}}$ be its image under~$\sigma$, twisted at the vertices of the parity of $v_0$ and standard at the
others.

Restrict $D_{\mathrm{tw}}$ to $\Gamma_B=\Gamma^+(\lambda_0)$, which acts freely on edges. Choose an $\F_{p^2}$-basis of the
fibre at one edge of each $\Gamma_B$-orbit and transport it by $\Gamma_B$. Every fibre becomes $\F_{p^2}$, and evaluation of
sections at the special points turns the datum into a $\Gamma_B$-equivariant local datum over $\F_{p^2}$ in the sense of
\cref{sec:prelim-tanner}. Its local codes are the codes of \cref{lem:grs} with $q=p^2$: evaluation is injective because
$m\le p^2-1$, so the Tanner code of the geometric datum on $\Gamma\backslash\Tree$ is the Tanner code $T_{\mathrm{tw}}(\Gamma\backslash\Tree)$
of \cref{sec:prelim-tanner} for every $\Gamma\subseteq\Gamma_B$ of finite index, and the local checks interpolate on
$t=p^2-m$ coordinates. Extending scalars to $\bar\F_p$ gives the Tanner code of $D_{\mathrm{tw}}\otimes\bar\F_p$, which is
therefore $T_{\mathrm{tw}}(\Gamma\backslash\Tree)\otimes_{\F_{p^2}}\bar\F_p$.

\begin{lemma}[The geometric datum is twisted]\label{lem:geo}
\notformalized
The datum $D_{\mathrm{geo}}$ extends to a $G$-equivariant geometric datum of degree $m=k(p-1)$, under the identification
$\Psi$ of $U_0$ with $\Q_{p^2}^2$ of \cref{lem:unif}, whose action extends that of $\Gamma^+$; its type is standard at the
vertices of one parity and twisted at the others.
\end{lemma}
\begin{informalproof}
The group $J$ consists of the automorphisms of the isocrystal that commute with $F$, respect the $\OO_E$-grading and
multiply the polarization by a factor in $\Q_p^\times$~\cite[(4.5)]{KudlaRapoport1999}. Such an automorphism is determined
by its action on $U_0$, because it commutes with $\Phi$, and any $g\in\GL(U_0)$ with $\det g\in\Q_p^\times$ extends to
one. So conjugation by $\Psi$ identifies $J$ with $\{g\in\GL_2(\Q_{p^2}):\det g\in\Q_p^\times\}$, as
in~\cite[(5.12)]{KudlaRapoport1999}; if $\Psi$ is $\sigma$-semilinear, conjugation by it applies $\sigma$ to determinants,
which preserves $\Q_p^\times$. In particular $G\subseteq J^0$. The group $J$
acts on $\mathcal M^{(n)}_{\mathrm{red}}$ and on $L_k$ compatibly with the crossings and with the identification of the dual
graph with $\Tree$~\cite[Section~6.3]{Luck2026}, so $D_{\mathrm{geo}}$ is $G$-equivariant, and the action of $\Gamma^+$ of
\cref{lem:unif} is the restriction.

Let $\Lambda$ be a lattice of the first kind and $K$ its stabilizer in $G$. A point of $C_\Lambda$ is a Dieudonn\'e
lattice $M=M_0\oplus M_1$ with $M_0=\Lambda\otimes W$, for which the index $0$ is critical, $F^2M_0=pM_0$, equivalently
$V^2M_0=pM_0$~\cite[Definition~4.1, Lemma~4.3 and~(4.9)]{KudlaRapoport1999}. The group $K$ acts on
$C_\Lambda=\Pone(\Lambda/p\Lambda)$ through its action on $\Lambda/p\Lambda$~\cite[Lemma~4.3]{KudlaRapoport1999}. After a
change of coordinates this is the action through $\rho_i$ if $\Psi$ is linear, and through $\sigma\circ\rho_i$ if $\Psi$
is $\sigma$-semilinear, where $v_i$ is the vertex among $v_0,v_1$ of the parity of $[\Psi\Lambda]$; in the second case
$\Psi$ induces a $\sigma$-semilinear bijection of $\Lambda/p\Lambda$ onto $\Psi\Lambda/p\Psi\Lambda$. If $g\in K$ is congruent to $1$ modulo $p\Lambda$, then $g$ fixes $C_\Lambda$ pointwise. It
also acts trivially on $VM/pM$, the fibre of the Hodge bundle at the point $M$ in the contravariant Dieudonn\'e theory
of~\cite{KudlaRapoport1999}. Indeed $(g-1)M_0\subseteq pM_0$, and, since $g$ commutes with $F$ and $V$
and $FM_1\subseteq M_0$~\cite[(4.6)]{KudlaRapoport1999}, $(g-1)M_1=F^{-1}(g-1)FM_1\subseteq F^{-1}pM_0=VM_0$. Hence
$(g-1)VM_1=V(g-1)M_1\subseteq V^2M_0=pM_0$ and $(g-1)VM_0=V(g-1)M_0\subseteq pVM_0\subseteq pM_1$. So $g$ acts
trivially on the fibres of $\omega_1$, $\omega_2$ and $L_k$ over $C_\Lambda$, and the stabilizer acts on $(C_\Lambda,L_k)$
through its reduction modulo~$p$. For lines of the second kind the same holds with the roles of $M_0$ and $M_1$ exchanged,
and the action on $C_{\Lambda'}$, in the coordinates of $U_0$ transported by $\Phi$, is conjugated by $\sigma$ because
$\Phi$ is $\sigma$-semilinear and commutes with $J$~\cite[Lemma~6.9, steps (1)--(2)]{Luck2026}. So if $\Psi$ is linear,
the type is standard at the vertices that carry lines of the first kind and twisted at the others; if $\Psi$ is
$\sigma$-semilinear, both are replaced by their images under $\sigma$, as in~\cite[Lemma~6.9, step~(5)]{Luck2026}. Lines
of the first kind sit at the vertices of one parity and lines of the second kind at the others, so in both cases the type
is standard at one parity and twisted at the other. By \cref{lem:unif}(c) the degree is~$m$.
\end{informalproof}

\begin{lemma}[Rigidity]\label{lem:rigid}
\notformalized
Two $G$-equivariant geometric data of degree $m$ with the same types at $v_0$ and at $v_1$ are isomorphic,
$G$-equivariantly. If moreover $-1\in G$ acts trivially on both data, then for every subgroup
$\Gamma\subseteq\PSL_2(\Q_{p^2})$ acting freely on $\Tree$ with finite quotient, their Tanner codes on $\Gamma\backslash\Tree$
are isomorphic.
\end{lemma}
\begin{informalproof}
Let $D$ and $D'$ be the data. (1) At $v_i$ the types agree, so there is a $K_i$-equivariant isomorphism
$\psi_i\colon C_{v_i}\to C'_{v_i}$. (2) The bundles $L_{v_i}$ and $\psi_i^*L'_{v_i}$ have the same degree, hence are
isomorphic; the two actions of $K_i$ on $L_{v_i}$ that result factor through $\SL_2(\F_{p^2})$ and cover the same action on
the curve, so they differ by a character of $\SL_2(\F_{p^2})$ with values in $\bar\F_p^\times$. Such a character is trivial:
$\SL_2(\F_{p^2})$ is generated by the elementary matrices $(\begin{smallmatrix}1&a\\0&1\end{smallmatrix})$ and
$(\begin{smallmatrix}1&0\\a&1\end{smallmatrix})$, which have order dividing $p$, and $\bar\F_p^\times$ has no element of
order~$p$. So the isomorphism of bundles can be chosen $K_i$-equivariant. (3) The special
point $P_{v_0,e_0}$ is fixed by the Iwahori subgroup, whose image in the action on $C_{v_0}$ is a Borel subgroup of
$\SL_2(\F_{p^2})$; it contains a nontrivial unipotent element, which has exactly one fixed point on $\Pone$. So
$\psi_0(P_{v_0,e_0})=P'_{v_0,e_0}$, and likewise at~$v_1$. (4) The two identifications of fibres at $e_0$ now differ by a
scalar; rescaling the isomorphism of bundles at $v_1$ makes them agree. (5) Transport the isomorphisms by $G$, setting
$\psi_{gv_i}=g\psi_ig^{-1}$; this is well defined because $\psi_i$ is $K_i$-equivariant. Every edge is $ge_0$ for some
$g\in G$, so the special points and identifications match at every edge, and the transported maps form a $G$-equivariant
isomorphism of data. It descends to $\Gamma\backslash\Tree$ and maps one Tanner code onto the other.
\end{informalproof}

\begin{remark}[The twist is forced]\label{rem:forced}
For $k\ne(p+1)/2$, equivariance alone forces the twist, independently of the computation with lines of the second kind
in the proof of \cref{lem:geo}: a $G$-equivariant geometric datum of degree $m$ whose types at $v_0$ and $v_1$ are each
standard or twisted is standard at one of them and twisted at the other. In a datum that is standard at
both $v_0$ and $v_1$, the Iwahori subgroup acts on the fibres at the two special points of $e_0$ by $\bar a^{-m}$ and
$\bar d^{-m}=\bar a^m$, and an equivariant identification exists only if $\bar a^{2m}=1$ for all $\bar a\in\F_{p^2}^\times$,
that is, if $p+1$ divides $2k$. For $2\le k\le p-2$ this means $k=(p+1)/2$. The same holds when both types are twisted.
The weight $k=(p+1)/2$ of the twelfth paper, where $p=3$ and $k=2$, is exactly the exceptional case, which is why that paper
needed an explicit identification of the type~\cite[Lemma~6.9]{Luck2026}.
\end{remark}

\begin{theorem}[Theorem G]\label{thm:G}
\notformalized
Let $p\ge7$ be a prime inert in $E=\Q(\sqrt5)$, let $2\le k\le p-2$, $m=k(p-1)$ and $t=p^2-m$, and let $\lambda_0$ be a
prime of $\OO_E$ of prime norm $\ell\ge7$. The twisted datum $D_{\mathrm{tw}}$ is a $\Gamma^+(\lambda_0)$-equivariant local
datum over $\F_{p^2}$ on the $(p^2+1)$-regular tree, whose local codes are generalized Reed--Solomon codes of length
$p^2+1$ and dimension $m+1$, with local checks interpolating on $t$ coordinates. For every $j\ge1$,
\[
  \dim_{\F_{p^2}}T_{\mathrm{tw}}(X_j)\ \ge\ (k-1)^2c_j=\frac{(k-1)^2}{p^2+1}\,\abs{X_j^1}.
\]
\end{theorem}
\begin{informalproof}
The first sentence summarizes the construction of the twisted model. By \cref{lem:geo}, $D_{\mathrm{geo}}$ is a
$G$-equivariant datum of degree $m$, standard at the vertices of one parity and twisted at the others. If the vertices
of the parity of $v_0$ are standard, \cref{lem:rigid} gives $D_{\mathrm{geo}}\cong D_{\mathrm{tw}}\otimes\bar\F_p$;
otherwise $D_{\mathrm{geo}}\cong D^\sigma_{\mathrm{tw}}\otimes\bar\F_p$, whose Tanner codes are the images under $\sigma$
of those of $D_{\mathrm{tw}}$. In both cases, by \cref{lem:unif}(b), \cref{lem:sections,lem:RRk},
\[
  \dim_{\F_{p^2}}T_{\mathrm{tw}}(X_j)=\dim_{\bar\F_p}T_{D_{\mathrm{geo}}}(X_j)=h^0(Y,L_k|_Y)\ge(k-1)^2c_j.
\]
Finally $\abs{X_j^1}=(p^2+1)c_j$ by \cref{lem:group}(b).
\end{informalproof}

Counting constraints gives only $\dim T_{\mathrm{tw}}(X_j)\ge(p^2+1-2t)c_j$, since each vertex imposes $t$ conditions. For
$t$ near $(2-\sqrt3)(p^2+1)$ this is less than half the length, while Theorem G gives a rate of at least about
$(\sqrt3-1)^2\approx0.536$. The excess comes from the geometry of the Goren--Oort configuration and is a phenomenon of
characteristic $p$ (\cref{sec:limits}).

\begin{remark}[A numerical check]\label{rem:numericalG}
The script \path{tools/twisted_tanner.py} builds the twisted datum for $p=7$ and a prime $\lambda_0$ above $\ell=11$,
where $c_1=11$, and computes $\dim_{\F_{49}}T_{\mathrm{tw}}(X_1)$ exactly. For $k=2,3,4,5$ it finds $22$, $44$, $111$
and~$176$, against the bounds $11$, $44$, $99$ and~$176$ of Theorem~G. As controls, the untwisted gluing, which is
equivariant only for $k=(p+1)/2=4$ (\cref{rem:forced}), gives $60<99$ there, and random gauges on the edges give the
counting bound $\max(0,(p^2+1-2t)c_1)$ or one dimension more. Here $\ell<4p$, so $X_1$ is not simple, which Theorem~G
does not need. The computation is not used in any proof.
\end{remark}
